\section{方法}

为在受限机载算力与高昂观测代价下做出稳健的航路选择，本文提出如图\ref{fig:overview}(d)所示的四步闭环决策方法：首先通过按需精算消除数值歧义，其次在基于多年海洋大数据标定的合理海流集合内利用投影梯度上升搜索最坏合理反例，最后根据反例的最坏后悔执行门控分流。本节依次阐述各模块的设计原理、其所替代的常规做法及完整算法流程。

\subsection{按需精算：数值歧义的快速识别与消除}
区分求解器误差与环境不确定性的首要原则，是仅当粗网格求解的账面时间差落在数值离散误差容限之内时，才动用高分辨率网格进行精算。常规的确定性航路规划往往采取一刀切的做法：要么为了节省算力对所有任务一律使用粗算，导致在两侧航时接近时因网格锯齿效应而选错方向；要么不计代价对所有任务一律执行耗时64倍的精算，在那些粗图上两侧原本就相差十几个小时的明朗任务上白白浪费机载电能。这种无差别的算力投放忽略了数值离散误差具有明确统计上界这一事实。

为此，我们在开发年份的710个任务上预先离线对比粗算（8千米网格、8方向）与精算（2千米网格、32方向）在同一张粗图 $\boldsymbol{u}_0$ 上求出的两侧时间差 $\Delta T_{\text{coarse}}(\boldsymbol{u}_0)$ 与 $\Delta T_{\text{fine}}(\boldsymbol{u}_0)$。统计二者绝对偏差 $|\Delta T_{\text{coarse}}(\boldsymbol{u}_0) - \Delta T_{\text{fine}}(\boldsymbol{u}_0)|$ 的经验分布，取其95\%分位数作为数值误差容限 $\tau_{\text{num}}$（在开发年份标定得 $\tau_{\text{num}} = 3.2$ 小时）。在线执行时，航行器首先以相对计算量为1的粗算快速求得 $\Delta T_{\text{coarse}}(\boldsymbol{u}_0)$；若 $|\Delta T_{\text{coarse}}(\boldsymbol{u}_0)| \ge \tau_{\text{num}}$，说明精算有95\%以上的把握不会改变粗图上的左右快慢顺序，直接采纳粗算结果；仅当 $|\Delta T_{\text{coarse}}(\boldsymbol{u}_0)| < \tau_{\text{num}}$ 时，才判定存在数值歧义并触发精算。这一按需精算机制以极小的平均算力开销确保了地图基准选择 $s_0 = \arg\min_{s \in \{\text{L}, \text{R}\}} T_s(\boldsymbol{u}_0)$ 不受求解器粗糙度的误导。

\subsection{反例检验：基于包络定理的合理集合投影梯度搜索}
判断粗图推荐的一侧 $s_0$ 是否面临严重的环境歧义，最直接的准则是主动在合理集合 $\mathcal{U}(r)$ 内寻找一个使所选一侧 $s_0$ 相比另一侧 $\bar{s}_0$ 尽可能变慢的“最坏合理反例”，而非在少数随机样本上求平均。现有的集合路径规划方法通常从误差分布中随机抽取 $M$ 个（如32个）流场样本，统计各样本下的航时均值或翻转比例\cite{wang2016path}。然而，翻转比例只反映有多少样本翻转了决定，不反映翻转后损失有多大；而两侧时间差只对少数扰动方向敏感，随机样本未必落在这些方向上。

为直接揪出这一最坏流场，我们将反例检验构建为合理球域上的约束最大化问题。定义任意流场 $\boldsymbol{u}(\boldsymbol{z}) \in \mathcal{U}(r)$ 下选用一侧 $s$ 的“最坏后悔”函数为该侧航时超出另一侧 $\bar{s}$ 航时的部分，并在半径 $r$ 的合理球内求其最大值 $R_s(r)$：
\begin{equation}
R_s(r) = \max_{\|\boldsymbol{z}\|_2 \le r} J_s(\boldsymbol{z}), \qquad J_s(\boldsymbol{z}) = \max\left\{ 0,\, T_s(\boldsymbol{u}(\boldsymbol{z})) - T_{\bar{s}}(\boldsymbol{u}(\boldsymbol{z})) \right\}.
\label{eq:worst_regret}
\end{equation}
当针对地图所选一侧 $s_0$ 求出的 $R_{s_0}(r) > 0$ 时，即表明合理集合内存在使另一侧 $\bar{s}_0$ 反而更快的合理反例，且 $R_{s_0}(r)$ 给出了若真实海洋恰为该反例流场时坚持选 $s_0$ 将多付出的最大航时（小时）。

求解式(\ref{eq:worst_regret})的关键难点在于目标函数 $T_{\text{L}}(\boldsymbol{u}(\boldsymbol{z})) - T_{\text{R}}(\boldsymbol{u}(\boldsymbol{z}))$ 嵌套了图最短路算子，若对每个模态系数 $z_k$ 采用有限差分求导，每次迭代将需要 $2K$ 次最短路求解，机载算力根本无法承受。本文利用最优化中的包络定理（Danskin定理）避免了这一开销：当流场发生微小扰动时，最短路径 $P_s^*(\boldsymbol{z})$ 自身的节点序列在一阶近似下保持不变，因此最优航时 $T_s(\boldsymbol{u}(\boldsymbol{z}))$ 对流场切变系数 $\boldsymbol{z}$ 的梯度，精确等于沿着当前已求出的那条最快航线 $P_s^*(\boldsymbol{z})$ 将各条边通行时间 $t(e)$ 对局部流速的偏导数直接累加：
\begin{equation}
\frac{\partial T_s}{\partial z_k} = \sum_{e \in P_s^*(\boldsymbol{z})} \left( \frac{\partial t(e)}{\partial u_{\parallel}} \boldsymbol{d}_e + \frac{\partial t(e)}{\partial u_{\perp}} \boldsymbol{n}_e \right) \cdot \frac{\boldsymbol{m}_k(\boldsymbol{x}_e^{-}) + \boldsymbol{m}_k(\boldsymbol{x}_e^{+})}{2},
\label{eq:envelope_grad}
\end{equation}
其中 $\boldsymbol{x}_e^{-}$、$\boldsymbol{x}_e^{+}$ 为有向边 $e$ 的两端节点，$\boldsymbol{n}_e$ 为垂直边方向的单位向量，$\boldsymbol{m}_k$ 为式(\ref{eq:plausible_field})中的误差模态，而标量导数由式(\ref{eq:edge_time})直接解析给出：
\begin{equation}
\frac{\partial t(e)}{\partial u_{\parallel}} = -\frac{\ell_e}{v_{\text{g}}(e)^2}, \qquad \frac{\partial t(e)}{\partial u_{\perp}} = \frac{\ell_e u_{\perp}}{v_{\text{g}}(e)^2 \sqrt{V^2 - u_{\perp}^2}}.
\label{eq:edge_deriv}
\end{equation}
得益于式(\ref{eq:envelope_grad})，每一步梯度计算只需在标准网格（4千米、16方向）上对左、右两侧各求一次最短路（共2次求解），即可同时获得两侧航时差对全部 $K$ 个流函数模态系数的完整梯度向量 $\boldsymbol{g}(\boldsymbol{z}) = \nabla_{\boldsymbol{z}} \Delta T(\boldsymbol{u}(\boldsymbol{z}))$。

在具体搜索策略上，若目标函数 $\Delta T(\boldsymbol{u}(\boldsymbol{z}))$ 关于 $\boldsymbol{z}$ 完全线性，则球约束下的极值必然精确落在从原点 $\boldsymbol{z}=\boldsymbol{0}$ 沿初始梯度 $\boldsymbol{g}(\boldsymbol{0})$ 直接跳到球面边界的点上。利用这一物理特性，我们的反例搜索先在粗图中心 $\boldsymbol{z}=\boldsymbol{0}$ 处计算初始梯度，一步直接跳至半径为 $r$ 的球面边界 $\boldsymbol{z}^{(1)} = \pm r\, \boldsymbol{g}(\boldsymbol{0}) / \|\boldsymbol{g}(\boldsymbol{0})\|_2$，随后在球内执行步长减半的投影梯度上升以适应流线弯曲带来的非线性效应。为同时评估左、右两侧各自的最坏后悔并支持不同半径下的快速查询，算法分别朝使左侧变慢（最大化 $\Delta T$）与使右侧变慢（最小化 $\Delta T$）两个方向、在基准半径 $r=1$ 与 $r=2$ 上各搜索一次，其他任意标定半径 $r$ 下的最坏后悔直接由 $r=1$ 与 $r=2$ 的结果线性插值得到。整个反例检验过程在每个任务上共计消耗约71次标准网格航线求解，与评估一个32成员集合所需的计算量（$32 \times 2 = 64$次求解）基本持平。

\subsection{海洋大数据标定：合理集合的物理结构与覆盖半径}
“合理反例”之所以具有实际决策指导意义而非脱离现实的数学构造，关键在于“合理”二字严格由多年历史海洋大数据标定，而非人为主观指定。若像一般的鲁棒决策那样凭经验设定不确定性集合的形状与半径\cite{benhaim2006infogap}，在海洋中要么因半径过小而漏掉黑潮锋面的真实摆动，要么因在弱流区也施加过大扰动而制造出不现实的反例。

本文利用开发年份（2019--2021年，365张快照）中3天前1°粗图 $\boldsymbol{u}_0$ 减去当天1/12°真实海流 $\boldsymbol{u}_{\text{true}}$ 的历史误差样本，分两步完成合理集合的物理标定。第一步标定空间结构：统计研究海区各网格点上历史误差流速的大小，生成空间非均匀的典型误差幅度图（即式(\ref{eq:plausible_field})中缩放矩阵 $\mathbf{S}(\boldsymbol{x})$ 的依据；3天前1°粗图在全区的典型误差中位数为0.17米/秒，沿黑潮主轴显著偏大，而当天0.5°图仅为0.07米/秒），并计算两点间误差相关系数随空间距离的衰减曲线，定出3天前1°粗图的误差相关长度 $L_{\text{corr}} = 76$ 千米（当天0.5°图为18千米），以此生成流函数模态 $\{\phi_k\}$。

第二步标定合理集合的半径 $r$：什么样的半径才算既不过于悲观、又能真实反映海洋可能发生的变化？我们提出基于决策目标的覆盖率准则——在开发年份的710个任务上，考察随半径 $r$ 增大，合理集合 $\mathcal{U}(r)$ 所给出的两侧时间差范围 $[\min_{\boldsymbol{u}\in\mathcal{U}(r)}\Delta T(\boldsymbol{u}),\, \max_{\boldsymbol{u}\in\mathcal{U}(r)}\Delta T(\boldsymbol{u})]$ 能够包住当天真实海流两侧时间差 $\Delta T(\boldsymbol{u}_{\text{true}})$ 的任务比例。将覆盖比例达到至少90\%的最小半径选定为标定半径，由此在开发年份上唯一定出 $r = 1.5$（该半径在开发年份实际覆盖93\%的任务，在完全独立的测试年份覆盖94\%的任务）。这一准则以决策关心的两侧时间差为尺度，使合理集合恰好大到能包住大多数真实海洋，而不必凭经验猜测半径。

\subsection{门控响应：依据最坏后悔与观测代价权衡行动}
在获得地图所选一侧的最坏后悔 $R_{s_0}(r)$ 与另一侧的最坏后悔 $R_{\bar{s}_0}(r)$ 后，最后一步是根据最坏后悔是否超过门控阈值 $\theta$ 以及外部观测代价的高低，在“照地图走”“上浮取新图”与“改走保守一侧”之间做出响应。现有的自适应采样或保守规划方法往往走两个极端：要么只要检测到不确定性就不顾代价地去测量，忽略了原地停留、绕行或频繁上浮本身就会直接吃掉数小时航时；要么对所有任务无条件采用保守或贝叶斯期望策略，在地图本来就十分可靠的低风险任务上错失最快航线。

本文设计的门控逻辑遵循“小风险直接走、大风险比代价”的原则。设开发年份上按平均损失最小选定的门控阈值为 $\theta$（标定得 $\theta = 8$ 小时）：
\begin{enumerate}
    \item 若 $R_{s_0}(r) \le \theta$，说明即使合理集合内最不利的流场成真，坚持按地图选 $s_0$ 至多也只会遭受不超过 $\theta$ 的有限劣势，环境歧义尚不足以推翻地图，航行器直接照地图走（选择 $s_0$），无需耗费额外时间观测。
    \item 若 $R_{s_0}(r) > \theta$，说明当前决策暴露在严重的尾部风险之下。此时进一步比较可用观测手段的净收益与不观测情况下的保守选择：当存在代价低廉且能有效消解大尺度流场歧义的观测手段（如上浮接收当天0.5°图，且单次上浮代价 $c_{\text{surf}}$ 较低）时，触发“门控后上浮取图”，在更新后的新图上重新选边出发；反之，当上浮通信受海况或任务隐蔽性限制而代价高昂，或仅有无法看清远方流场的原位测流可用时，触发“门控后取保守”，直接无需观测地改选合理集合内最坏后悔更小的一侧：
    \begin{equation}
    s_{\text{cons}} = \arg\min_{s \in \{\text{L}, \text{R}\}} R_s(r).
    \label{eq:conservative_choice}
    \end{equation}
\end{enumerate}
算法\ref{alg:counterexample_gate}给出了整个合理反例检验与门控决策的完整伪代码流程。

\begin{algorithm}[htbp]
\caption{基于合理反例检验的自主水下航行器航路门控决策算法}
\label{alg:counterexample_gate}
\begin{algorithmic}[1]
\Require 粗分辨率海流图 $\boldsymbol{u}_0$，标定误差幅度图 $\sigma(\boldsymbol{x})$ 与无散流函数基 $\{\phi_k\}_{k=1}^{K}$，标定半径 $r=1.5$，数值误差容限 $\tau_{\text{num}}=3.2$ 小时，门控阈值 $\theta=8.0$ 小时，上浮取图代价 $c_{\text{surf}}$ 及代价平衡点 $c^* \approx 1.0$ 小时。
\Ensure 最终执行的航路决策动作（选定绕行一侧 $s^*$ 或触发上浮取新图）。
\State 在粗图 $\boldsymbol{u}_0$ 上用粗算（8千米网格、8方向）求解两侧航时，得时间差 $\Delta T_{\text{coarse}}(\boldsymbol{u}_0)$。
\If{$|\Delta T_{\text{coarse}}(\boldsymbol{u}_0)| < \tau_{\text{num}}$}
    \State 触发按需精算：在 $\boldsymbol{u}_0$ 上用精算（2千米网格、32方向）重算得 $\Delta T_{\text{fine}}(\boldsymbol{u}_0)$，定地图倾向侧 $s_0 \gets \arg\min_{s \in \{\text{L},\text{R}\}} T_s^{\text{fine}}(\boldsymbol{u}_0)$。
\Else
    \State 直接采纳粗算结果：$s_0 \gets \arg\min_{s \in \{\text{L},\text{R}\}} T_s^{\text{coarse}}(\boldsymbol{u}_0)$。
\EndIf
\State 在中心 $\boldsymbol{z}=\boldsymbol{0}$ 处用标准网格（4千米、16方向）求左、右最快航线，按包络定理式(\ref{eq:envelope_grad})计算初始梯度 $\boldsymbol{g}(\boldsymbol{0})$。
\State 分别沿 $+\boldsymbol{g}(\boldsymbol{0})$ 与 $-\boldsymbol{g}(\boldsymbol{0})$ 方向跳至 $r_1=1$ 与 $r_2=2$ 的球面边界，执行步长减半的投影梯度上升，并在标定半径 $r$ 处线性插值得两侧最坏后悔 $R_{\text{L}}(r)$ 与 $R_{\text{R}}(r)$。
\If{$R_{s_0}(r) \le \theta$}
    \State \Return 直接照地图走：选择绕行一侧 $s^* \gets s_0$。
\Else
    \If{允许上浮通信且 $c_{\text{surf}} < c^*$}
        \State \Return 门控后上浮取图：付出代价 $c_{\text{surf}}$ 获取当天0.5°海流图 $\boldsymbol{u}_{\text{new}}$，选 $s^* \gets \arg\min_{s \in \{\text{L},\text{R}\}} T_s(\boldsymbol{u}_{\text{new}})$。
    \Else
        \State \Return 门控后取保守：选择最坏后悔更小的一侧 $s^* \gets \arg\min_{s \in \{\text{L},\text{R}\}} R_s(r)$。
    \EndIf
\EndIf
\end{algorithmic}
\end{algorithm}
